% ============================================================================
%  ESTADÍSTICA INFERENCIAL — Sesión 9
%  Semana 9: 5 al 10 de octubre de 2026
%  2.4 Distribución t de Student
%  2.5 Distribución de las medias muestrales y Teorema del Límite Central
% ============================================================================
\documentclass[aspectratio=169,10pt]{beamer}

\usepackage[utf8]{inputenc}
\IfFileExists{spanish.ldf}{\usepackage[spanish,es-noshorthands]{babel}}{}
\usepackage{amsmath,amssymb}
\usepackage{booktabs}
\usepackage{tikz}
\usepackage{listings}

\usetheme{Uniminuto}

\uniasignatura{Estadística inferencial}
\uniprofesor{Andrés Felipe Puerta Velez}
\unicorreo{andres.puerta-v@uniminuto.edu.co}
\uninrc{CBASBL061}
\unisemana{Sesión 9 — Semana 9 · 5 al 10 de octubre de 2026}

\newcommand{\listatiny}{\setbeamerfont{itemize/enumerate body}{size=\tiny}\tiny}

\newcommand{\fx}[1]{{\ttfamily\color{uniAzul}#1}}

% ------------------------------------------------------------- CÓDIGO R ----
\lstdefinestyle{Rcode}{
  language=R,
  basicstyle=\ttfamily\scriptsize,
  keywordstyle=\color{uniAzul}\bfseries,
  commentstyle=\color{uniGris}\itshape,
  stringstyle=\color{uniTitulo},
  showstringspaces=false, numbers=none, breaklines=true,
  columns=fullflexible,
  backgroundcolor=\color{uniCremaClaro},
  frame=leftline, framerule=1.2pt, rulecolor=\color{uniNaranja},
  xleftmargin=0.35cm, framexleftmargin=0.35cm,
  framexrightmargin=0.15cm, framextopmargin=2pt, framexbottommargin=2pt,
  aboveskip=0.15cm, belowskip=0.15cm
}
\lstdefinestyle{Rmini}{
  language=R,
  basicstyle=\ttfamily\tiny,
  keywordstyle=\color{uniAzul}\bfseries,
  commentstyle=\color{uniGris}\itshape,
  stringstyle=\color{uniTitulo},
  showstringspaces=false, numbers=none, breaklines=true,
  columns=fullflexible,
  backgroundcolor=\color{uniCremaClaro},
  frame=leftline, framerule=1.2pt, rulecolor=\color{uniNaranja},
  xleftmargin=0.35cm, framexleftmargin=0.35cm,
  framexrightmargin=0.15cm, framextopmargin=2pt, framexbottommargin=2pt,
  aboveskip=0.12cm, belowskip=0.12cm
}
\lstdefinestyle{Rout}{
  basicstyle=\ttfamily\tiny\color{black},
  showstringspaces=false, numbers=none, breaklines=true,
  columns=fullflexible,
  backgroundcolor=\color{uniGris!10},
  frame=leftline, framerule=1.2pt, rulecolor=\color{uniGris},
  xleftmargin=0.35cm, framexleftmargin=0.35cm,
  framexrightmargin=0.15cm, framextopmargin=2pt, framexbottommargin=2pt,
  aboveskip=0.12cm, belowskip=0.12cm
}

\begin{document}

\portada

\begin{frame}{Orden del día}
  \begin{enumerate}
    \item Llamado a lista y revisión de los ejercicios propuestos
    \item Repaso: dónde quedamos
    \item \azul{Tema 2.5: distribución de las medias muestrales y el TLC}
    \item \azul{Tema 2.4: distribución $t$ de Student}
    \item Ejercicios en clase
    \item Cierre y ejercicios propuestos
  \end{enumerate}
\end{frame}

\begin{frame}{Repaso: dónde quedamos}
  Ya sabemos calcular probabilidades para \azul{un} dato: si
  $X\sim N(500,50)$, ¿qué probabilidad hay de que \emph{un} pedido tarde más de
  600 minutos?

  \vspace{0.25cm}
  \begin{alertblock}{La pregunta de hoy, que es la del curso entero}
    \footnotesize
    Ya no preguntamos por un dato sino por el \destacar{promedio de una
    muestra}. Si tomo 36 pedidos, ¿qué probabilidad hay de que su promedio
    $\bar{x}$ supere los 830 minutos? Esta es la pregunta que permite pasar de
    la muestra a la población.
  \end{alertblock}

  \vspace{0.15cm}
  {\footnotesize\color{uniGris} La respuesta se llama \azul{distribución
  muestral de la media} y es la bisagra entre la probabilidad y la
  inferencia\,[1].}
\end{frame}

\seccion{2.5 Distribución de las medias muestrales}

\begin{frame}{La media muestral es una variable aleatoria}
  Si de una misma población se toman muchas muestras de tamaño $n$, cada una da
  un $\bar{x}$ distinto. Esos valores tienen su propia distribución.

  \vspace{0.25cm}
  \begin{block}{Media y error estándar de $\bar{X}$}
    \[ E(\bar{X})=\mu \qquad\qquad
       \sigma_{\bar{X}}=\frac{\sigma}{\sqrt{n}} \]
  \end{block}

  \vspace{0.2cm}
  \begin{alertblock}{Lo que dicen estas dos fórmulas}
    \footnotesize
    \begin{itemize}
      \item El promedio de los promedios da en el blanco: $\bar{X}$ es un
            estimador \azul{insesgado} de $\mu$.
      \item La dispersión de $\bar{X}$ es \destacar{menor} que la de $X$, y se
            reduce con $\sqrt{n}$. Promediar suaviza.
    \end{itemize}
  \end{alertblock}
\end{frame}

\begin{frame}{Teorema del Límite Central (TLC)}
  \begin{block}{Enunciado}
    \small
    Si el tamaño de muestra $n$ es suficientemente grande, la distribución de
    $\bar{X}$ es \destacar{aproximadamente normal}
    \[ \bar{X}\;\sim\;N\!\left(\mu,\;\frac{\sigma}{\sqrt{n}}\right) \]
    \azul{sin importar cómo sea la distribución de la población original}.
  \end{block}

  \vspace{0.2cm}
  \begin{columns}[T,onlytextwidth]
    \begin{column}{0.48\textwidth}
      \begin{exampleblock}{Regla práctica}
        \footnotesize
        $n\ge 30$ basta en la mayoría de los casos. Si la población ya es
        normal, vale para cualquier $n$.
      \end{exampleblock}
    \end{column}
    \begin{column}{0.48\textwidth}
      \begin{alertblock}{Por qué es tan importante}
        \footnotesize
        Permite usar la tabla normal aunque no sepamos nada de la forma de la
        población. Es lo que hace posible todo lo que viene: estimación,
        intervalos y pruebas de hipótesis.
      \end{alertblock}
    \end{column}
  \end{columns}
\end{frame}

\begin{frame}{Estandarización de la media muestral}
  \begin{block}{Un dato frente a un promedio}
    \[ \text{Un dato: } Z=\frac{X-\mu}{\sigma}
       \qquad\qquad
       \text{Un promedio: } Z=\frac{\bar{X}-\mu}{\sigma/\sqrt{n}} \]
  \end{block}

  \vspace{0.25cm}
  \begin{alertblock}{El error más frecuente del curso}
    \footnotesize
    Olvidar el $\sqrt{n}$ del denominador. Si la pregunta habla de
    \destacar{un promedio, una media o un total de $n$ observaciones}, va
    $\sigma/\sqrt{n}$. Si habla de \destacar{un solo elemento}, va $\sigma$.
  \end{alertblock}

  \vspace{0.2cm}
  {\footnotesize\color{uniGris} Con $\sigma=120$ y $n=36$, el error estándar es
  $120/6=20$: la dispersión de los promedios es seis veces menor que la de los
  datos individuales.}
\end{frame}

\begin{frame}{Ejercicio 1 — en clase}
  \begin{alertblock}{Individual (9 minutos)}
    El gasto mensual de los clientes de una tienda tiene
    \destacar{$\mu=800$} y \destacar{$\sigma=120$} (miles de pesos). Se toma
    una muestra aleatoria de \destacar{$n=36$} clientes.

    \vspace{0.25cm}
    \begin{enumerate}
      \item ¿Cuál es el error estándar de $\bar{X}$?
      \item $P(\bar{X}>830)$
      \item $P(770<\bar{X}<830)$
      \item $P(\bar{X}<780)$
    \end{enumerate}
  \end{alertblock}
\end{frame}

\begin{frame}{Ejercicio 1 — solución}
  \small
  \begin{block}{Error estándar y estandarización}
    \[ \sigma_{\bar{X}}=\frac{120}{\sqrt{36}}=\azul{20}
       \qquad Z=\frac{\bar{x}-800}{20} \]
    \begin{align*}
      \bar{x}=830 &\Rightarrow z=1{,}50 &&P(\bar{X}>830)=1-0{,}9332=\azul{0{,}0668}\\
      \bar{x}=770 &\Rightarrow z=-1{,}50 &&P(770<\bar{X}<830)=\azul{0{,}8664}\\
      \bar{x}=780 &\Rightarrow z=-1{,}00 &&P(\bar{X}<780)=\azul{0{,}1587}
    \end{align*}
  \end{block}

  \vspace{0.1cm}
  {\footnotesize\color{uniGris} Compare: para \emph{un} cliente,
  $P(X>830)$ daría $z=0{,}25$ y una probabilidad de $0{,}40$. El promedio de 36
  clientes es mucho más estable que un cliente suelto.}
\end{frame}

\begin{frame}{Distribución muestral de la proporción}
  \begin{block}{El mismo razonamiento con proporciones}
    \[ E(\hat{p})=p \qquad
       \sigma_{\hat p}=\sqrt{\frac{p(1-p)}{n}} \qquad
       Z=\frac{\hat{p}-p}{\sqrt{p(1-p)/n}} \]
  \end{block}

  \vspace{0.2cm}
  \begin{exampleblock}{Ejemplo}
    \footnotesize
    Si el 40\% de los clientes usa la app ($p=0{,}40$) y se toma una muestra
    de $n=100$:
    \[ \sigma_{\hat p}=\sqrt{\tfrac{0{,}4\cdot 0{,}6}{100}}=0{,}0490
       \qquad P(\hat{p}>0{,}45)=\azul{0{,}1537} \]
  \end{exampleblock}

  \vspace{0.1cm}
  {\scriptsize\color{uniGris} La aproximación normal vale si $np\ge 5$ y
  $n(1-p)\ge 5$. Aquí $np=40$ y $n(1-p)=60$: sin problema.}
\end{frame}

\seccion{2.4 Distribución t de Student}

\begin{frame}{El problema: casi nunca conocemos $\sigma$}
  Para calcular $Z=(\bar{X}-\mu)/(\sigma/\sqrt{n})$ hace falta $\sigma$, un
  parámetro poblacional. En la práctica no se conoce y hay que reemplazarlo por
  la desviación muestral $s$.

  \vspace{0.25cm}
  \begin{block}{Lo que resulta ya no es normal}
    \[ T=\frac{\bar{X}-\mu}{s/\sqrt{n}} \;\sim\; t_{\,n-1} \]
    Distribución $t$ de Student con \destacar{$n-1$ grados de libertad}.
  \end{block}

  \vspace{0.2cm}
  {\footnotesize\color{uniGris} Al estimar $\sigma$ con $s$ se introduce
  incertidumbre adicional, y la distribución se vuelve más ancha en las colas
  para compensarla.}
\end{frame}

\begin{frame}{Características de la $t$}
  \begin{columns}[T,onlytextwidth]
    \begin{column}{0.48\textwidth}
      \begin{block}{Se parece a la normal, pero}
        \footnotesize
        \begin{itemize}
          \item Es simétrica y centrada en 0.
          \item Tiene \azul{colas más pesadas}: valores extremos son más
                probables.
          \item Depende de los grados de libertad $gl=n-1$.
          \item A medida que $gl$ crece se acerca a la normal; con $gl>30$ la
                diferencia es mínima.
        \end{itemize}
      \end{block}
    \end{column}
    \begin{column}{0.48\textwidth}
      \begin{exampleblock}{Valores críticos bilaterales al 95\%}
        \footnotesize
        \begin{center}
        \begin{tabular}{@{}cr@{}}
          \toprule
          $gl$ & $t$ \\
          \midrule
          9   & 2{,}262 \\
          14  & 2{,}145 \\
          19  & 2{,}093 \\
          29  & 2{,}045 \\
          $\infty$ & 1{,}960 \\
          \bottomrule
        \end{tabular}
        \end{center}
      \end{exampleblock}
    \end{column}
  \end{columns}

  \vspace{0.15cm}
  {\scriptsize\color{uniGris} Note cómo el valor crítico baja hacia 1{,}96: la
  $t$ converge a la normal.}
\end{frame}

\begin{frame}{¿Cuándo $Z$ y cuándo $t$?}
  \footnotesize
  \begin{center}
  \begin{tabular}{@{}llll@{}}
    \toprule
    \textbf{Situación} & \textbf{$\sigma$} & \textbf{$n$} & \textbf{Se usa} \\
    \midrule
    Población normal & Conocida   & cualquiera & $Z$ \\
    Población normal & Desconocida & cualquiera & $t$ con $gl=n-1$ \\
    Población cualquiera & Conocida & $n\ge 30$ & $Z$ (por TLC) \\
    Población cualquiera & Desconocida & $n\ge 30$ & $t$; en la práctica
      $\approx Z$ \\
    Población cualquiera & Desconocida & $n<30$ & $t$, y hay que suponer
      normalidad \\
    \bottomrule
  \end{tabular}
  \end{center}

  \vspace{0.25cm}
  \begin{alertblock}{Regla de bolsillo}
    \footnotesize
    \destacar{Si en el enunciado aparece $s$ (desviación de la muestra), use
    $t$. Si aparece $\sigma$ (de la población), use $Z$.}
  \end{alertblock}
\end{frame}

\begin{frame}{Ejercicio 2 — en clase}
  \begin{alertblock}{En parejas (9 minutos)}
    Para cada situación indique si se usa $Z$ o $t$, y con cuántos grados de
    libertad cuando corresponda.

    \vspace{0.25cm}
    \begin{enumerate}
      \item $n=40$, población no normal, $\sigma=15$ conocida.
      \item $n=12$, población normal, $s=4{,}2$.
      \item $n=200$, $s=30$.
      \item $n=8$, población normal, $\sigma=5$ conocida.
      \item $n=25$, población normal, $s=6$.
    \end{enumerate}
  \end{alertblock}
\end{frame}

\begin{frame}{Ejercicio 2 — solución}
  \small
  \begin{block}{Respuestas}
    \footnotesize
    \begin{enumerate}
      \item \azul{$Z$}: $\sigma$ conocida y $n\ge 30$, el TLC garantiza
            normalidad de $\bar{X}$.
      \item \azul{$t$ con $gl=11$}: $\sigma$ desconocida y población normal.
      \item \azul{$t$ con $gl=199$}, prácticamente igual a $Z$ porque
            $gl$ es enorme.
      \item \azul{$Z$}: aunque $n$ es pequeño, la población es normal y
            $\sigma$ se conoce.
      \item \azul{$t$ con $gl=24$}, valor crítico bilateral al 95\%:
            $2{,}064$.
    \end{enumerate}
  \end{block}

  \vspace{0.1cm}
  {\footnotesize\color{uniGris} Lo decisivo es siempre lo mismo: ¿el dato de
  dispersión es $\sigma$ o es $s$?}
\end{frame}

\begin{frame}{Ejercicio 3 — en clase}
  \begin{alertblock}{Individual (8 minutos)}
    En una empresa, el \destacar{40\%} de los clientes usa la aplicación
    móvil. Se toma una muestra de \destacar{$n=100$} clientes.

    \vspace{0.25cm}
    \begin{enumerate}
      \item Calcule $\sigma_{\hat p}$.
      \item Verifique que se puede usar la aproximación normal.
      \item $P(\hat{p}>0{,}45)$
      \item $P(0{,}35<\hat{p}<0{,}45)$
    \end{enumerate}
  \end{alertblock}
\end{frame}

\begin{frame}{Ejercicio 3 — solución}
  \begin{columns}[T,onlytextwidth]
    \begin{column}{0.52\textwidth}
      \begin{block}{Cálculos}
        \footnotesize
        \[ \sigma_{\hat p}=\sqrt{\frac{0{,}40(0{,}60)}{100}}
           =\sqrt{0{,}0024}=\azul{0{,}0490} \]
        \begin{align*}
          \hat p=0{,}45 &\Rightarrow z=\tfrac{0{,}05}{0{,}0490}=1{,}02\\
          &\Rightarrow P(\hat p>0{,}45)=\azul{0{,}1537}\\[0.2em]
          P(0{,}35<\hat p<0{,}45)&=\azul{0{,}6926}
        \end{align*}
      \end{block}
    \end{column}
    \begin{column}{0.44\textwidth}
      \begin{exampleblock}{Punto 2}
        \footnotesize
        $np=100(0{,}40)=40\ge 5$ y $n(1-p)=60\ge 5$. La aproximación normal es
        válida.
      \end{exampleblock}
      \begin{alertblock}{Lectura}
        \footnotesize
        Cerca del 69\% de las muestras de 100 clientes darán una proporción
        entre 35\% y 45\%. Ese margen de $\pm 5$ puntos es, en esencia, un
        intervalo de confianza: el tema de la próxima sesión.
      \end{alertblock}
    \end{column}
  \end{columns}
\end{frame}

\begin{frame}[fragile]{La $t$ y el TLC, en R}
  \scriptsize
\begin{lstlisting}[style=Rmini]
# Valores criticos de t (reemplazan la tabla)
qt(0.975, df = 19)      # bilateral 95%, gl = 19
qt(0.975, df = 24)      # bilateral 95%, gl = 24
qt(0.975, df = 199)     # con gl grande se parece a 1.96
pt(-2.5, df = 15)       # area a la izquierda de t = -2.5

# El TLC, visto en vivo: 5000 medias de muestras de tamano 30
set.seed(2026)
poblacion <- rexp(100000, rate = 1/50)      # poblacion MUY asimetrica
medias    <- replicate(5000, mean(sample(poblacion, 30)))
hist(medias, main = "Distribucion de la media muestral", col = "steelblue")
c(mean(medias), sd(medias), 50/sqrt(30))    # comparar con mu y sigma/raiz(n)
\end{lstlisting}
\begin{lstlisting}[style=Rout]
[1] 2.093024
[1] 2.063899
[1] 1.971957
[1] 0.01222196
\end{lstlisting}

  \vspace{-0.12cm}
  \begin{alertblock}{Por qué vale la pena correr las cuatro últimas líneas}
    \tiny
    La población es exponencial —muy asimétrica— y aun así el histograma de las
    5\,000 medias sale \destacar{acampanado}, con desviación cercana a
    $\sigma/\sqrt{n}$. Eso es el Teorema del Límite Central dejando de ser un
    enunciado y volviéndose una imagen.
  \end{alertblock}
\end{frame}

\begin{frame}{Cierre de la sesión}
  \begin{block}{Lo que debe quedar claro hoy}
    \begin{itemize}
      \item $\bar{X}$ es aleatoria: $E(\bar X)=\mu$ y
            $\sigma_{\bar X}=\sigma/\sqrt{n}$.
      \item TLC: con $n\ge 30$, $\bar{X}$ es aproximadamente normal sea cual
            sea la población.
      \item Si se usa $s$ en vez de $\sigma$, la distribución es $t$ con
            $n-1$ grados de libertad.
      \item Nunca olvide el $\sqrt{n}$ cuando la pregunta sea sobre un
            promedio.
    \end{itemize}
  \end{block}

  \vspace{0.2cm}
  \begin{alertblock}{Trabajo independiente para la próxima sesión}
    \footnotesize
    Resolver los \destacar{20 ejercicios de la lámina siguiente}. En la sesión
    10 empezamos a \azul{estimar}: estimación puntual e intervalos de
    confianza.
  \end{alertblock}
\end{frame}

\begin{frame}{Ejercicios propuestos}
  \vspace{-0.1cm}
  \begin{beamercolorbox}[wd=\textwidth,sep=0.10cm]{block body}
    \centering\scriptsize
    \textbf{1--10:} $\mu=800$, $\sigma=120$, $n=36$. \quad
    \textbf{11--14:} $p=0{,}40$, $n=100$.
  \end{beamercolorbox}
  \listatiny
  \vspace{0.15cm}
  \begin{columns}[T,onlytextwidth]
    \begin{column}{0.47\textwidth}
      \begin{enumerate}\setlength{\itemsep}{0.16cm}
        \item Error estándar $\sigma_{\bar X}$
        \item $z$ para $\bar x=830$
        \item $P(\bar X>830)$
        \item $P(\bar X<780)$
        \item $P(770<\bar X<830)$
        \item $P(\bar X>800)$
        \item Error estándar si $n=144$
        \item $P(\bar X>830)$ con $n=144$
        \item ¿Por cuánto hay que multiplicar $n$ para reducir el error
              estándar a la mitad?
        \item $P(X>830)$ para \emph{un} cliente
      \end{enumerate}
    \end{column}
    \begin{column}{0.50\textwidth}
      \begin{enumerate}\setlength{\itemsep}{0.16cm}
        \setcounter{enumi}{10}
        \item $\sigma_{\hat p}$
        \item ¿Se cumple $np\ge 5$ y $n(1-p)\ge 5$?
        \item $P(\hat p>0{,}45)$
        \item $P(0{,}35<\hat p<0{,}45)$
        \item ¿$Z$ o $t$? $n=40$, $\sigma$ conocida
        \item ¿$Z$ o $t$? $n=12$, $s$ conocida
        \item ¿$Z$ o $t$? $n=200$, $s$ conocida
        \item $gl$ si $n=25$
        \item $t$ bilateral 95\% con $gl=14$
        \item $t$ bilateral 95\% con $gl=29$
      \end{enumerate}
    \end{column}
  \end{columns}
\end{frame}

\begin{frame}{Respuestas}{Autoevaluación de los ejercicios propuestos}
  \listatiny
  \vspace{0.15cm}
  \begin{columns}[T,onlytextwidth]
    \begin{column}{0.47\textwidth}
      \begin{enumerate}\setlength{\itemsep}{0.16cm}
        \item $20$
        \item $1{,}50$
        \item $0{,}0668$
        \item $0{,}1587$
        \item $0{,}8664$
        \item $0{,}5000$
        \item $10$
        \item $0{,}0013$ \;{\color{uniGris}$z=3{,}0$}
        \item Por $4$
        \item $0{,}4013$ \;{\color{uniTitulo}$z=0{,}25$; sin $\sqrt{n}$}
      \end{enumerate}
    \end{column}
    \begin{column}{0.50\textwidth}
      \begin{enumerate}\setlength{\itemsep}{0.16cm}
        \setcounter{enumi}{10}
        \item $0{,}0490$
        \item Sí: $40$ y $60$
        \item $0{,}1537$
        \item $0{,}6926$
        \item $Z$
        \item $t$ con $gl=11$
        \item $t$ con $gl=199\approx Z$
        \item $gl=24$
        \item $2{,}145$
        \item $2{,}045$
      \end{enumerate}
    \end{column}
  \end{columns}
\end{frame}

\begin{frame}{Referencias}
  \footnotesize
  \begin{itemize}
    \item[{[1]}] Díaz Rodríguez, M. (2022). \textit{Estadística inferencial
      aplicada}. Universidad del Norte.
      \texttt{https://elibro.net/es/lc/uniminuto/titulos/221614}
    \item[{[2]}] Llinás Solano, H. (2014). \textit{Introducción a la
      estadística con aplicaciones en ciencias sociales}. Universidad del
      Norte. \texttt{https://elibro.net/es/lc/uniminuto/titulos/70062}
    \item[{[3]}] Contento Rubio, M. R. (2019). \textit{Estadística con
      aplicaciones en R}. Editorial Utadeo.
      \texttt{https://elibro.net/es/lc/uniminuto/titulos/220926}
  \end{itemize}

  \vspace{0.35cm}
  \begin{center}
    {\Large\color{uniAzul}\textbf{¿Preguntas?}}
  \end{center}
\end{frame}

\end{document}
